\section*{第 \ref{ch:g11:prob} 章 --- 概率与随机变量}

\begin{solution}{exo:g11:prob:1}
$52$ 张牌上的古典概型。红心：$\frac{13}{52} = \frac14$。国王：
$\frac{4}{52} = \frac{1}{13}$。红心或国王：
$\frac{13 + 4 - 1}{52} = \frac{16}{52} = \frac{4}{13}$（红心国王
只计一次）。都不是：$1 - \frac{4}{13} = \frac{9}{13}$。
\end{solution}

\begin{solution}{exo:g11:prob:2}
$S = 6$：数对 $(1,5), (2,4), (3,3), (4,2), (5,1)$，于是
$\P(S = 6) = \frac{5}{36}$。

$S \geq 10$：和为 $10$ 的有 3 个数对，和为 $11$ 的有 2 个，和为 $12$ 的有 1 个：
$\P(S \geq 10) = \frac{6}{36} = \frac16$。

$S$ 为偶数时，和为 $2, 4, 6, 8, 10, 12$ 的数对个数依次为 $1 + 3 + 5 + 5 + 3 + 1 = 18$，
即共 18 个数对，概率为 $\frac{18}{36} = \frac12$。
\end{solution}

\begin{solution}{exo:g11:prob:3}
各概率之和为 $1$：$p = 1 - 0.3 - 0.2 - 0.4 = 0.1$。于是
\[
\E(X) = 0.3 \times (-1) + 0.2 \times 0 + 0.4 \times 2 + 0.1 \times 5
= -0.3 + 0.8 + 0.5 = 1 .
\]
\end{solution}

\begin{solution}{exo:g11:prob:4}
$\E(Y) = 2 \times 3 - 5 = 1$; $\V(Y) = 2^2 \times 4 = 16$;
$\sigma(Y) = 4$.
\end{solution}

\begin{solution}{exo:g11:prob:5}
奖金分别为 $0$、$2$、$5$ 欧元，概率依次为 $0.5$、$0.3$、
$0.2$；扣除 $2$ 欧元赌注后，$G$ 的取值为 $-2, 0, 3$：
\begin{center}
\begin{tabular}{c|ccc}
$g$ & $-2$ & $0$ & $3$\\
\hline
$\P(G=g)$ & $0.5$ & $0.3$ & $0.2$
\end{tabular}
\end{center}
$\E(G) = -1 + 0 + 0.6 = -0.4$：玩家每局平均
损失 $40$ 欧分 --- 对玩家不利。
\end{solution}

\begin{solution}{exo:g11:prob:6}
依次摸球：$\P(X = 2) = \frac35 \times \frac24 = \frac{3}{10}$
（先红后红）；$\P(X = 0) = \frac25 \times \frac14 = \frac{1}{10}$
（先蓝后蓝）。由此
$\P(X = 1) = 1 - \frac{3}{10} - \frac{1}{10} = \frac{6}{10}$。于是
\[
\E(X) = 0 \times \frac1{10} + 1 \times \frac6{10} + 2 \times \frac3{10}
= \frac{12}{10} = 1.2 ,
\]
与直觉相符：平均而言，两次摸球中红球所占的比例为 $\frac35$。
\end{solution}

\begin{solution}{exo:g11:prob:7}
$\E(G^2) = 0.5 \times 4 + 0.3 \times 0 + 0.2 \times 9 = 3.8$，由
简便公式可得
$\V(G) = 3.8 - (-0.4)^2 = 3.8 - 0.16 = 3.64$，且
$\sigma(G) \approx 1.9$ 欧元。
\end{solution}

\begin{solution}{exo:g11:prob:8}
收益为 $80$ 减去赔付额：
\begin{center}
\begin{tabular}{c|ccc}
$b$ & $80$ & $-920$ & $-4920$\\
\hline
$\P(B=b)$ & $0.94$ & $0.05$ & $0.01$
\end{tabular}
\end{center}
$\E(B) = 75.2 - 46 - 49.2 = -20$：按此价格，公司平均每份保单\emph{亏损}
$20$ 欧元。保费相对于所承保的风险过低
（保费须高于 $100$ 欧元，
才能使期望收益为正）。
\end{solution}

\begin{solution}{exo:g11:prob:9}
随机得分为 $+1$ 的概率为 $\frac14$，
为 $x$ 的概率为 $\frac34$：
\[
\E = \frac14 + \frac{3x}{4} = 0 \iff x = -\frac13 .
\]
当罚分为 $-\frac13$ 时，随机猜测平均而言没有收益。
\end{solution}

\begin{solution}{exo:g11:prob:10}
两枚硬币：结果相同（HH 或 TT）的概率为 $\frac24 \times 2 =
\frac12$。玩家的收益为 $+1$ 或 $-1$，概率相等：
两位玩家的期望均为 $0$ --- 游戏公平。

三枚硬币：三枚结果全部相同（HHH 或 TTT）的概率为 $\frac{2}{8} =
\frac14$。玩家获得 $2$ 的概率为 $\frac14$，
付出 $1$ 的概率为 $\frac34$：
$\E = \frac24 - \frac34 = -\frac14$。每轮游戏对安妮有利
$25$ 欧分；若安妮为三枚全同的结果支付 $3$ 欧元，则游戏公平。
\end{solution}

\begin{solution}{exo:g11:prob:11}
\emph{1.} 组织者收取 $1000 \times 5 = 5000$ 欧元，
共支付 $2000 + 5 \times 200 + 20 \times 25 = 3500$ 欧元，利润
为\emph{定值} $1500$（所有票售出后即无随机性）。
单个玩家的收益 $G$（奖金减去 $5$ 欧元的票价）：
\begin{center}
\begin{tabular}{c|cccc}
$g$ & $1995$ & $195$ & $20$ & $-5$\\
\hline\rule{0pt}{11pt}%
$\P(G=g)$ & $\dfrac{1}{1000}$ & $\dfrac{5}{1000}$ & $\dfrac{20}{1000}$
& $\dfrac{974}{1000}$
\end{tabular}
\end{center}

\emph{2.}
\[
\E(G) = \frac{1995 + 5 \times 195 + 20 \times 20 - 974 \times 5}{1000}
= \frac{1995 + 975 + 400 - 4870}{1000} = -1.5 .
\]
每名玩家期望损失 $1.50$ 欧元；$1000$ 名玩家合计
损失 $1500$ 欧元，恰为组织者的利润。账目是平衡的：
玩家平均的损失，恰是
组织者所收的钱。
\end{solution}

\begin{solution}{pb:g11:prob:1}
\textbf{1.} $G$ 取值 $-2, -1, 0, 1, 2, 3$，各值的概率均为
$\frac16$，于是 $\E(G) = \frac{-2 - 1 + 0 + 1 + 2 +
3}{6} = \frac12$ 欧元。

\textbf{2.} $\E(G^2) = \frac{4 + 1 + 0 + 1 + 4 + 9}{6} =
\frac{19}{6}$; $V(G) = \frac{19}{6} - \frac14 = \frac{35}{12}
\approx 2.92$; $\sigma(G) \approx 1.71$.

\textbf{3.} $\E(2G - 1) = 2 \times \frac12 - 1 = 0$;
$V(2G - 1) = 4\,V(G) = \frac{35}{3}$.

\textbf{4.} $\E(\text{和}) = 7$：公平价格为 $7$ 欧元。当价格定为 $8$ 时，
期望损失为每局 $1$ 欧元。

\textbf{5.} 公平保费：$0.05 \times 800 = 40$ 欧元。
收取 $60$ 欧元时，每份保单的期望利润为 $20$ 欧元。

\textbf{6.} 按 $5:3$ 分配，奖励的是\emph{过去} --- 但奖池
应归属于赢得\emph{未来}的人，而 B 仍
有可能获胜。全部给 A，则是把可能的胜利
当作必然，B 的小概率并非一文不值。

\textbf{7.} 还需进行的局数为 $3$；每局结果（A 胜/B 胜）
的全部情形为 AAA、AAB、ABA、ABB、BAA、BAB、BBA、BBB。玩家 B 只有
三局全胜（即仅 BBB 一种）才能赢得整场。A 赢得整场的
情形有 $7$ 种（共 $8$ 种），公平分配为
$64 \times \frac78 = 56$ 皮斯托尔归 A，$8$ 皮斯托尔归 B。

\textbf{8.} 多打的无用局数不改变
任何人的胜负：A 一旦获得第六胜，之后各局
便只是走过场。但将局数固定为恰好 $3$ 局，
使 $8$ 种结果等可能 --- 诚实的计数要求
情节等长，而把它们补足毫无代价。

\textbf{9.} 比分 $5$--$5$ 时，还需一局决胜，各得 $32$。比分
为 $5$--$4$ 时，A 得 $\frac{64 + 32}{2} = 48$。比分为 $5$--$3$ 时，
下一局的结果：A 胜则为 $64$，否则进入 $5$--$4$ 局面，其价值为
$48$：A 的份额为 $\frac{64 + 48}{2} = 56$ ---
费马的 $56 : 8$，由此倒推即可重新推出。

\textbf{10.} A 还需 $2$ 局，B 还需 $3$ 局：假想再赛 $4$ 局
（共 $16$ 种结果）。B 胜 $3$ 次或 $4$ 次即赢得整场，
即 $4 + 1 = 5$ 种结果。分配为：
$B$: $64 \times \frac{5}{16} = 20$; $A$: $44$.

\textbf{11.} 对不确定的未来收益而言，其公平现值即为它的期望，
即按概率加权的各种可能结局的平均值。保险（第 5 题）
与金融（为不确定收益定价）
都是这句话演化而成的行业；
赌场亦然，只不过站在了桌子的另一边。

\textbf{12.} 奖池终究要分完：两人的份额
是随机变量，其和为常数 $64$，因此由
线性性得 $\E_A + \E_B = 64$ --- 公平的份额恰好分完奖池，
无论比分如何。

\textbf{13.} 由线性性，$\E(X + Y) = 3.5 + 3.5 = 7$，无需列表，
即使两颗粘在一起的骰子也成立。单颗骰子的方差为：
$\E(X^2) - 3.5^2 = \frac{91}{6} - \frac{49}{4} = \frac{35}{12}$;
若两颗骰子相互独立，
$V(X + Y) = \frac{35}{12} + \frac{35}{12} = \frac{35}{6}$.

\textbf{14.} 每位固定的客人取回自己帽子的概率为
$\frac1n$（其帽子是 $n$ 顶中的一顶）。将
这 $n$ 位客人求和：
$\E(X) = n \times \frac1n = 1$ --- 平均有一位幸运的客人，
无论宴会上有两顶帽子还是两千顶。

\textbf{15.} 当 $n = 2$ 时，分配方式为两人全对（$X = 2$）或
两人互换（$X = 0$）：$\E = 1$。当 $n = 3$ 时，六种分配给出
$X = 3, 1, 1, 1, 0, 0$: $\E = \frac{6}{6} = 1$.

\textbf{16.} 上述计算中，把 $n$ 个
单独的指示变量（“客人 $i$ 拿回了自己的帽子”）的期望相加，而
线性性对期望求和时\emph{无条件}成立 ---
相依性会破坏乘积或方差，却从不破坏
期望之和。这正是线性性的超能力：
无需过问独立性。

\textbf{17.} 两种情形的收益：$\E = 0.5 \times 10 - 5 = 0$ 及
$0.005 \times 1000 - 5 = 0$。收益的方差：
A：$\E(G^2) = 0.5 \times 25 + 0.5 \times 25 = 25$；
B：$0.005 \times 995^2 + 0.995 \times 25 \approx 4\,975$：
$\sigma_A = 5$，$\sigma_B \approx 70.5$。人们购买 B，
他们所买的并非期望（两者均为零），而是
\emph{离散程度} --- 一个小概率的、能改变人生的
右尾。

\textbf{18.} 第 $n$ 次抛掷才首次出现反面的概率为
$2^{-n}$，奖金为 $2^n$：每个 $n$ 对期望的贡献为
$2^{-n} \times 2^n = 1$，因此期望为
$1 + 1 + 1 + \dots$：无穷大。然而没有理智的人会愿意
花 $1\,000$ 欧元玩一次，因为以概率 $\frac{15}{16}$，
奖金至多为 $16$。期望概括的是长期
\emph{可重复}赌局的结果；一次性、尾部极端的游戏
不能以其均值来定价。

\textbf{19.} 期望：投保为 $-20$；不投保为
$-0.01 \times 1000 = -10$。平均而言，不投保的策略更优
--- 对于能承受损失、并多次重复该赌局的人而言。
无法承受 $-1\,000$ 欧元损失的家庭，并不属于
“取平均”的范畴：它愿意多付 $10$ 欧元，以消除
方差。保险公司取平均，家庭则购买确定性。

\textbf{20.} 它诞生于两位急不可耐的赌徒分配奖池，
借由线性性获得力量，它把希望相加，
无须征询相依性的许可；对风险视而不见，而方差
--- 彩票、圣彼得堡游戏与这家庭的保单
--- 本应由它来度量；并于次年由大数
定律加冕，它解释了为何平均值总会
最终“兑现”。
\end{solution}
